\documentclass[11pt]{article}
\usepackage[a4paper,margin=25mm]{geometry}
\usepackage{amsmath,amssymb,amsthm,hyperref}
\newtheorem{theorem}{Theorem}
\title{A trace-norm vertex that is not a partial permutation}
\author{Ziang Ni (\texttt{ziangni-sys})}
\date{10 October 2026}
\begin{document}
\maketitle
\noindent\textbf{Conjecture 00000007151.}
The source asserts that the extremal body of the trace norm in
Schatten $p$ has vertices at partial permutation matrices.
With the usual norm-unit-body and unsigned partial-permutation
conventions, this is false already for real $1\times1$ matrices.
Neither language imposes entrywise positivity or a dimension restriction.

\paragraph{Actual matrix norm.}
For a real matrix $A$, its singular values are the nonnegative square
roots of the eigenvalues of $A^T A$; its trace norm is their sum.
Take the one-dimensional real matrix space $M_1(\mathbb R)$.
For $A=[a]$, actual matrix multiplication gives $A^TA=[a^2]$.
Its sole eigenvalue is $a^2$: multiplication on a vector $[v]$ is
$[a^2v]$, and a nonzero eigenvector forces the eigenvalue to be $a^2$.
Thus its sole singular value is $\sqrt{a^2}=|a|$, and
\[
 \|[a]\|_1=|a|,\qquad
 B=\{A\in M_1(\mathbb R):\|A\|_1\le1\}
   =\{[a]:-1\le a\le1\}.
\]
This is the actual trace-norm unit body. It is a line segment,
so the word ``vertex'' is unambiguous here. Choosing $p=1$ already
suffices; in fact every one-dimensional Schatten norm has the same
unit body.

\begin{theorem}
The matrix $V=[-1]$ is an exposed vertex of $B$, but is not a standard
partial permutation matrix.
\end{theorem}
\begin{proof}
Consider the nonzero matrix linear functional $\ell([a])=-a$.
For every $[a]\in B$, one has $\ell([a])\le1$, with equality
if and only if $a=-1$. Therefore its maximizing face on $B$ is
exactly $\{V\}$, so $V$ is exposed.
For completeness, suppose $V=tA+(1-t)C$ with $A,C\in B$ and
$0<t<1$. Writing $A=[a]$ and $C=[c]$, we get
\[
 0=t(a+1)+(1-t)(c+1).
\]
Both summands are nonnegative and their multipliers are positive;
hence $a=c=-1$. This also proves extremality directly.

A standard partial permutation matrix has entries in $\{0,1\}$,
with at most one $1$ in each row and each column. In dimension one
the only such matrices are $[0]$ and $[1]$. The entry $-1$ of $V$
is neither $0$ nor $1$, proving the claimed exclusion.
\end{proof}

\paragraph{Scope and formal verification.}
The example uses the usual unsigned partial-permutation definition,
not a modified definition allowing negative signs. The statement
contains no positivity condition that would remove $V$.
The Lean project uses actual \texttt{Matrix Unit Unit}\,$\mathbb R$,
actual transpose/multiplication and vector action, the uniquely
computed Gram eigenvalue, and its singular-value trace norm.
It proves the singleton maximizing face of an actual linear functional,
membership in mathlib's \texttt{extremePoints}, and failure of the
full zero-one row/column partial-permutation conditions.
The final existential theorem gives the counterexample.
Lean 4.19.0 and mathlib commit
\texttt{c44e0c8ee63ca166450922a373c7409c5d26b00b} are pinned;
the printed axiom audits contain only the permitted standard axioms.
\end{document}
